\begin{workedexample}[Worked Example: Patch resonant length]
A patch on FR-4 ($\varepsilon_r = 4.4$) at $f = 2.4\ \text{GHz}$ resonates
near a half wavelength in the substrate. Using
$\varepsilon_{\text{eff}} \approx (\varepsilon_r+1)/2 = 2.7$,
\[ L \approx \frac{c}{2f\sqrt{\varepsilon_{\text{eff}}}}
   = \frac{3\times10^{8}}{2(2.4\times10^{9})(1.64)} = 38\ \text{mm}, \]
trimmed slightly for fringing fields. Patches are low-profile and easily arrayed
but narrowband.
\end{workedexample}

\begin{keytakeaways}
\begin{itemize}[leftmargin=1.4em,itemsep=2pt]
\item Patch $L \approx c/(2f\sqrt{\varepsilon_{\text{eff}}})$: flat, cheap, arrayable, narrowband.
\item Aperture antennas (horn, reflector) give high gain $G = \eta_a 4\pi A/\lambda^2$.
\item Log-periodic/spiral trade size for bandwidth; PIFA/monopole suit handsets.
\end{itemize}
\end{keytakeaways}

\begin{exercises}
\item Patch length on a substrate with $\varepsilon_{\text{eff}} = 6.25$ at $1\ \text{GHz}$?
\item Which antenna family gives the widest bandwidth for its size?
\item Which gives the highest gain: a patch or a parabolic reflector?
\end{exercises}

\furtherreading{Balanis~\cite{balanis} (ch.~14); Pozar~\cite{pozar} (ch.~14).}
